\subsection{ソフトウェア工学経済の基礎}

\subsubsection{提案}

\begin{definitionbox}提案\par\smallskip
提案とは、検討対象となる一つの行動案である。ソフトウェア開発プロジェクトを実施する、既存部品を拡張する、別製品に置き換える、あるアルゴリズムを採用する、といった案が提案である。
\end{definitionbox}

提案は、実行するかしないかを判断する単位である。意思決定の目的は、組織の目標に最も合う提案を見つけることにある。提案は大規模プロジェクトだけではなく、設計、実装、テスト、運用の細かな選択にも現れる。

\subsubsection{キャッシュフロー}

キャッシュフローは、提案を実行した結果として組織に入る資金と出る資金の時系列である。ある月にサーバを購入すれば支出が発生し、製品の販売が始まれば収入が発生する。これらの個々の出入りをキャッシュフロー事象と呼び、一定期間にわたる全体の流れをキャッシュフロー列と呼ぶ。

\par\smallskip
\begin{center}
\centering
\IfFileExists{assets/generated_figures/ch15/fig-ch15-02.png}{\includegraphics[width=0.96\linewidth]{assets/generated_figures/ch15/fig-ch15-02.png}}{\fbox{\parbox[c][48mm][c]{0.9\linewidth}{\centering 画像生成待ち\\キャッシュフロー図}}}
\par\smallskip
\refstepcounter{figure}\label{fig:ch15-02}図 \thefigure: キャッシュフロー図
\end{center}
図 \thefigure は、キャッシュフロー図を視覚的に整理したものである。
\par\smallskip

\subsubsection{お金の時間価値}

お金には時間価値がある。同じ 100 万円でも、現在受け取る 100 万円と 3 年後に受け取る 100 万円は同じ価値ではない。現在の資金は投資でき、利息や収益を生む可能性があるためである。

基本的な現在価値の式は次のように表せる。
\[
PV = \frac{FV}{(1+i)^n}
\]
ここで、$PV$ は現在価値、$FV$ は将来価値、$i$ は利率、$n$ は期間である。初学者は、将来のお金を現在の価値へ割り引いて比較するための式だと理解すればよい。

\subsubsection{等価性}

等価性とは、異なる時点のキャッシュフローを同じ時点の価値に換算し、比較可能にする考え方である。時間価値を無視して「合計金額」だけを比べると、誤った判断になることがある。

\begin{notebox}{例}
案 A は今すぐ 100 万円の利益を生む。案 B は 5 年後に 110 万円の利益を生む。単純合計では案 B が大きく見えるが、現在価値へ換算すると案 A の方が有利な場合がある。
\end{notebox}

\subsubsection{比較基準}

比較基準とは、複数のキャッシュフロー列を同じ土俵で比べるための枠組みである。主な基準には次がある。

\par\smallskip\noindent\refstepcounter{table}\label{tab:ch15-03}表 \thetable: 比較基準の整理\par\smallskip
表 \thetable は、比較基準の整理を比較・確認しやすい形で整理したものである。\par\smallskip
\begin{longtable}{L{0.30\linewidth} L{0.65\linewidth}}
\toprule
比較基準 & 意味 \\
\midrule
現在価値 & すべての収入と支出を現在時点の価値に換算して比べる。 \\
将来価値 & すべての収入と支出を将来のある時点の価値に換算して比べる。 \\
年等価額 & 複数年の効果を、毎年同じ金額が発生するものとして比べる。 \\
\term{内部収益率}{Internal Rate of Return} & 提案の収益性を表す利率で比べる。 \\
割引回収期間 & 初期投資を、時間価値を考慮して何年で回収できるかを見る。 \\
\bottomrule
\end{longtable}

\subsubsection{代替案}

代替案とは、互いに比較して選択できる案の集合である。提案同士には依存関係や排他関係があることが多い。たとえば、提案 Y は提案 X を実行しなければ実行できないかもしれない。あるいは、提案 P と提案 Q は同時に実行できないかもしれない。

\term{意思決定}{Decision Making}では、提案の関係を整理し、互いに排他的な代替案へ変換することが重要である。特に、何もしない代替案も重要である。これは本当に何もしないという意味ではなく、「この候補群のどれも実行せず、別のことをする」という選択肢である。

\subsubsection{無形資産}

\begin{definitionbox}無形資産\par\smallskip
無形資産とは、物理的な形はないが、組織の成果や財務に影響する知識、経験、手順、文化、仕様、ツール、ノウハウなどである。知識資産とも呼ばれる。
\end{definitionbox}

ソフトウェアの提案は、無形資産に影響を受ける。たとえば、組織に高度な業務知識があるなら、その知識を活かす設計が価値を生む。逆に、現場の暗黙知を無視して導入したシステムは、技術的には正しくても業務に適合しない可能性がある。

\subsubsection{ビジネスモデル}

ビジネスモデルは、顧客は誰か、顧客は何に価値を感じるか、組織はどのように収益や持続可能性を得るか、どの費用構造で価値を届けるかを説明する枠組みである。

\begin{pointbox}{重要}
提案を評価するときは、ソフトウェア単体だけを見るのではなく、組織のビジネスモデルと無形資産を見る必要がある。これにより、隠れたリスクや機会を発見できる。
\end{pointbox}
